\documentclass{article}

\usepackage{a4}
\usepackage{url}
\usepackage{amsmath}
\usepackage{setspace}

\textwidth 6in
\oddsidemargin 0.15in
\topmargin -0.5in
\textheight 10in
\doublespacing

\begin{document}

\pagestyle{empty}
\thispagestyle{empty}

\begin{center}
    {\LARGE \bf \noindent KinecMamba: Anatomy-Aware 3D Human Pose Estimation using Mamba}
\end{center}

\vspace{0.2in}
{\bf \noindent Abstract}\\
\vspace{-0.1in}

\noindent Monocular 3D human pose estimation recovers the three-dimensional
positions of body joints from single-camera video. The task is usually solved in
two stages: a 2D detector locates joints on the image, and a lifting network maps
the resulting 2D keypoint sequence to 3D. Transformer-based lifting networks model
spatial and temporal dependencies well, but their self-attention scales
quadratically with the number of spatio-temporal tokens, which caps the temporal
context they can afford. State space models such as Mamba reach comparable modelling
quality at linear cost, yet they read their input as a single ordered sequence,
while a human skeleton has no natural order. Feeding the joints in raw index order
places anatomically distant joints next to each other and weakens the model's grasp
of local structure. This thesis presents KinecMamba, a bidirectional spatio-temporal
state space network for 2D-to-3D pose lifting. It preserves the linear cost of Mamba
through a four-directional cross-scan over the frame and joint axes, and orders the
joints along the spatial scan by a breadth-first traversal of the skeleton's
kinematic tree, from the pelvis outward along each limb. An anatomy-aware
bone-length loss supervises the lengths and temporal consistency of the
skeleton's bones, reinforcing the same structure the scan follows. KinecMamba
is evaluated on Human3.6M against Transformer- and state-space-based methods.
At roughly 1.26M parameters it reaches 41.55mm MPJPE and 34.52mm P-MPJPE,
improving on its closest comparable baseline, PoseMamba-S, on both metrics at
a similar parameter budget.

\end{document}

